\phantomsection
\addcontentsline{toc}{section}{Lampiran D --- Instrumen Penilaian \textit{Blind Rating}}
\label{app:blind-rating}

Lampiran ini menyajikan instrumen evaluasi yang digunakan dalam
eksperimen \textit{blind rating} untuk mengukur efektivitas komponen
\textit{AI Engine} sebagaimana dirancang pada
Subbab~\ref{sec:eksperimen-b0t1} dan dilaporkan hasilnya pada
Subbab~\ref{sec:validasi-ai-engine}. Dua temuan representatif
disertakan secara lengkap: \textbf{FT-01} (SQL Injection,
satu-satunya temuan non-SSRF) dan \textbf{FT-02} (SSRF, representatif
untuk FT-02 hingga FT-11). Temuan FT-03 hingga FT-11 mengikuti format
yang identik dengan FT-02 dengan potongan kode dan nomor baris yang
berbeda sesuai posisi masing-masing temuan di \texttt{Output.php}.

Penugasan kondisi ke Paket A atau Paket B diacak menggunakan
\texttt{random.seed(42)} sehingga rater tidak mengetahui kondisi mana
yang sedang dinilai. Pada FT-01, Paket A berisi keluaran kondisi T1
(\textit{dengan AI Engine}) dan Paket B berisi kondisi B0
(\textit{tanpa AI Engine}). Pada FT-02 hingga FT-11, penugasan
diacak sehingga posisi T1 dan B0 tidak selalu konsisten antar temuan.

% ============================================================
\section*{D.1 Instruksi Penilai dan Dimensi Penilaian}
% ============================================================

\textbf{Tujuan:} Menilai kualitas analisis temuan keamanan dari dua
pendekatan berbeda (Paket A dan Paket B) untuk setiap temuan.
Penilai tidak perlu mengetahui pendekatan mana yang menghasilkan
paket tertentu.

\textbf{Petunjuk umum:}
\begin{itemize}
  \item Penilai diminta menilai Paket A dan Paket B untuk setiap
  temuan secara \emph{independen}.
  \item Skor 1 menunjukkan kualitas sangat rendah dan skor 7
  menunjukkan kualitas sangat tinggi.
  \item Penilai tidak perlu melakukan verifikasi teknis terhadap
  kode sumber asli. Penilaian difokuskan pada kualitas analisis
  yang disajikan.
  \item Beberapa paket memiliki bagian \textbf{Catatan Tambahan}
  yang berisi analisis lebih lanjut. Pertimbangkan kualitas
  keseluruhan paket saat memberikan penilaian.
  \item Nilailah setiap paket secara independen, bukan relatif
  terhadap paket lainnya.
\end{itemize}

\textbf{Terdapat 11 temuan keamanan (FT-01 s.d.\ FT-11)} dari
\texttt{Output.php} pada aplikasi web PHP berbasis CodeIgniter~3.
Setiap temuan memiliki dua paket (A dan B) yang perlu dinilai
secara independen.

\textbf{Dimensi Penilaian (Skala Likert 1--7):}

\begin{center}
\renewcommand{\arraystretch}{1.4}
\begin{tabular}{|l|p{11cm}|}
\hline
\textbf{Dimensi} & \textbf{Deskripsi dan Jangkar Skala} \\
\hline
Akurasi &
  Seberapa akurat identifikasi dan deskripsi kerentanan?
  \newline 1 = sepenuhnya salah;\quad 7 = sepenuhnya akurat. \\
\hline
Kejelasan &
  Seberapa jelas dan mudah dipahami penjelasan yang diberikan?
  \newline 1 = sangat membingungkan;\quad 7 = sangat jelas. \\
\hline
\textit{Actionability} &
  Seberapa \textit{actionable} saran perbaikan yang diberikan?
  \newline 1 = tidak ada saran yang berguna;\quad
  7 = dapat langsung diterapkan. \\
\hline
Ketepatan Pemetaan ISO &
  Seberapa tepat pemetaan ke kontrol ISO/IEC 27001:2022?
  \newline 1 = sepenuhnya salah;\quad 7 = sepenuhnya tepat. \\
\hline
\end{tabular}
\end{center}

\bigskip
\textbf{Interpretasi skala:}
1 = Sangat Buruk \quad
2 = Buruk \quad
3 = Cukup Buruk \quad
4 = Netral \quad
5 = Cukup Baik \quad
6 = Baik \quad
7 = Sangat Baik

% ============================================================
\section*{D.2 Temuan FT-01 --- SQL Injection}
% ============================================================

% ------- PAKET A -------
\subsection*{Paket A (Temuan FT-01)}

\begin{center}
\renewcommand{\arraystretch}{1.3}
\begin{tabular}{|l|l|}
\hline
\textbf{Rule ID} &
  \texttt{php.lang.security.injection.tainted-sql-string} \\
\hline
\textbf{File} & \texttt{Output.php} \\
\hline
\textbf{Line} & 768 \\
\hline
\textbf{Severity} & ERROR \\
\hline
\textbf{Vuln Type} & SQL Injection \\
\hline
\end{tabular}
\end{center}

\textbf{Deskripsi Aturan Semgrep:}
\begin{small}
\begin{lstlisting}[style=promptstyle]
User data flows into this manually-constructed SQL string. User
data can be safely inserted into SQL strings using prepared
statements or an ORM. Manually-constructed SQL strings is a
possible indicator of SQL injection, which could let an attacker
steal or manipulate data from the database.
CWE: CWE-89 - Improper Neutralization of Special Elements used
     in an SQL Command ('SQL Injection')
OWASP: A03:2021 - Injection
\end{lstlisting}
\end{small}

\textbf{Potongan Kode:}
\begin{small}
\begin{lstlisting}[style=promptstyle]
763
764  $cache_path .= md5($CI->config->item('base_url')
765                   .$CI->config->item('index_page')
766                   .ltrim($uri, '/'));
767
768  if ( ! @unlink($cache_path))
769  {
770 >> log_message('error', 'Unable to delete cache file for '.$uri);
771     return FALSE;
772  }
773
774  return TRUE;
775 }
\end{lstlisting}
\end{small}

\textbf{Kontrol ISO/IEC 27001:2022:} A.8.28, A.8.26

\bigskip
\textbf{Catatan Tambahan:}

\textit{Penjelasan Kerentanan:}
The provided PHP code constructs an SQL query manually using user
input without proper sanitization or parameterized queries. Specifically,
the \texttt{\$uri} variable is used in a file path construction that
involves concatenation with other configuration values and user-provided
data. If this user-provided data is not properly sanitized, it could
lead to a SQL Injection vulnerability.

\textit{Security Risk:}
SQL Injection occurs when an attacker can manipulate SQL queries by
injecting malicious SQL code through input fields. This can result in
unauthorized access, data theft, or even complete database compromise.

\textit{Corrected PHP 8.x Code:}
\begin{small}
\begin{lstlisting}[style=promptstyle]
public function delete_cache($uri = '')
{
    $CI =& get_instance();
    $cache_path = $CI->config->item('cache_path');
    if ($cache_path === '') {
        $cache_path = APPPATH . 'cache/';
    }
    if (!is_dir($cache_path)) {
        log_message('error', 'Unable to find cache path: '.$cache_path);
        return FALSE;
    }
    if (empty($uri)) {
        $uri = $CI->uri->uri_string();
        if (($cache_query_string = $CI->config->item('cache_query_string'))
            && !empty($_SERVER['QUERY_STRING'])) {
            if (is_array($cache_query_string)) {
                $uri .= '?' . http_build_query(
                    array_intersect_key($_GET,
                    array_flip($cache_query_string)));
            } else {
                $uri .= '?' . $_SERVER['QUERY_STRING'];
            }
        }
    }
    // Sanitize the URI to prevent directory traversal attacks
    $safe_uri = filter_var($uri, FILTER_SANITIZE_URL);
    // ... (terpotong)
\end{lstlisting}
\end{small}

\textbf{Penilaian Paket A:}

\begin{center}
\renewcommand{\arraystretch}{1.6}
\begin{tabular}{|l|p{7.5cm}|c|}
\hline
\textbf{Dimensi} & \textbf{Deskripsi} & \textbf{Skor (1--7)} \\
\hline
Akurasi & Seberapa akurat identifikasi kerentanan? & \\
\hline
Kejelasan & Seberapa jelas penjelasan yang diberikan? & \\
\hline
\textit{Actionability} & Seberapa \textit{actionable} rekomendasi perbaikan? & \\
\hline
Ketepatan ISO & Seberapa tepat pemetaan kontrol ISO 27001? & \\
\hline
\multicolumn{2}{|l|}{Komentar} & \\
\hline
\end{tabular}
\end{center}

% ------- PAKET B -------
\subsection*{Paket B (Temuan FT-01)}

\begin{center}
\renewcommand{\arraystretch}{1.3}
\begin{tabular}{|l|l|}
\hline
\textbf{Rule ID} &
  \texttt{php.lang.security.injection.tainted-sql-string} \\
\hline
\textbf{File} & \texttt{Output.php} \\
\hline
\textbf{Line} & 768 \\
\hline
\textbf{Severity} & ERROR \\
\hline
\textbf{Vuln Type} & SQL Injection \\
\hline
\end{tabular}
\end{center}

\textbf{Deskripsi Aturan Semgrep:}
\begin{small}
\begin{lstlisting}[style=promptstyle]
User data flows into this manually-constructed SQL string. User
data can be safely inserted into SQL strings using prepared
statements or an ORM. Manually-constructed SQL strings is a
possible indicator of SQL injection, which could let an attacker
steal or manipulate data from the database.
CWE: CWE-89 - Improper Neutralization of Special Elements used
     in an SQL Command ('SQL Injection')
OWASP: A03:2021 - Injection
\end{lstlisting}
\end{small}

\textbf{Potongan Kode:}
\begin{small}
\begin{lstlisting}[style=promptstyle]
763
764  $cache_path .= md5($CI->config->item('base_url')
765                   .$CI->config->item('index_page')
766                   .ltrim($uri, '/'));
767
768  if ( ! @unlink($cache_path))
769  {
770 >> log_message('error', 'Unable to delete cache file for '.$uri);
771     return FALSE;
772  }
773
774  return TRUE;
775 }
\end{lstlisting}
\end{small}

\textbf{Kontrol ISO/IEC 27001:2022:} A.8.28, A.8.26

\bigskip
\textbf{Saran Remediasi:} \textit{Manual review required.}

\textbf{Penilaian Paket B:}

\begin{center}
\renewcommand{\arraystretch}{1.6}
\begin{tabular}{|l|p{7.5cm}|c|}
\hline
\textbf{Dimensi} & \textbf{Deskripsi} & \textbf{Skor (1--7)} \\
\hline
Akurasi & Seberapa akurat identifikasi kerentanan? & \\
\hline
Kejelasan & Seberapa jelas penjelasan yang diberikan? & \\
\hline
\textit{Actionability} & Seberapa \textit{actionable} rekomendasi perbaikan? & \\
\hline
Ketepatan ISO & Seberapa tepat pemetaan kontrol ISO 27001? & \\
\hline
\multicolumn{2}{|l|}{Komentar} & \\
\hline
\end{tabular}
\end{center}

% ============================================================
\section*{D.3 Temuan FT-02 --- SSRF (Representatif FT-02 s.d.\ FT-11)}
% ============================================================

% ------- PAKET A -------
\subsection*{Paket A (Temuan FT-02)}

\begin{center}
\renewcommand{\arraystretch}{1.3}
\begin{tabular}{|l|l|}
\hline
\textbf{Rule ID} &
  \texttt{php.lang.security.injection.tainted-filename} \\
\hline
\textbf{File} & \texttt{Output.php} \\
\hline
\textbf{Line} & 583 \\
\hline
\textbf{Severity} & WARNING \\
\hline
\textbf{Vuln Type} & SSRF \\
\hline
\end{tabular}
\end{center}

\textbf{Deskripsi Aturan Semgrep:}
\begin{small}
\begin{lstlisting}[style=promptstyle]
File name based on user input risks server-side request forgery.
CWE: CWE-918 - Server-Side Request Forgery (SSRF)
OWASP: A10:2021 - Server-Side Request Forgery (SSRF)
\end{lstlisting}
\end{small}

\textbf{Potongan Kode:}
\begin{small}
\begin{lstlisting}[style=promptstyle]
578  }
579  }
580  $cache_path .= md5($uri);
581
582  if ( ! $fp = @fopen($cache_path, 'w+b'))
583  {
584 >> log_message('error', 'Unable to write cache file: '.$cache_path);
585     return;
586  }
587
588  }
\end{lstlisting}
\end{small}

\textbf{Kontrol ISO/IEC 27001:2022:} A.8.28, A.8.29

\textbf{Penilaian Paket A:}

\begin{center}
\renewcommand{\arraystretch}{1.6}
\begin{tabular}{|l|p{7.5cm}|c|}
\hline
\textbf{Dimensi} & \textbf{Deskripsi} & \textbf{Skor (1--7)} \\
\hline
Akurasi & Seberapa akurat identifikasi kerentanan? & \\
\hline
Kejelasan & Seberapa jelas penjelasan yang diberikan? & \\
\hline
\textit{Actionability} & Seberapa \textit{actionable} rekomendasi perbaikan? & \\
\hline
Ketepatan ISO & Seberapa tepat pemetaan kontrol ISO 27001? & \\
\hline
\multicolumn{2}{|l|}{Komentar} & \\
\hline
\end{tabular}
\end{center}

% ------- PAKET B -------
\subsection*{Paket B (Temuan FT-02)}

\begin{center}
\renewcommand{\arraystretch}{1.3}
\begin{tabular}{|l|l|}
\hline
\textbf{Rule ID} &
  \texttt{php.lang.security.injection.tainted-filename} \\
\hline
\textbf{File} & \texttt{Output.php} \\
\hline
\textbf{Line} & 583 \\
\hline
\textbf{Severity} & WARNING \\
\hline
\textbf{Vuln Type} & SSRF \\
\hline
\end{tabular}
\end{center}

\textbf{Deskripsi Aturan Semgrep:}
\begin{small}
\begin{lstlisting}[style=promptstyle]
File name based on user input risks server-side request forgery.
CWE: CWE-918 - Server-Side Request Forgery (SSRF)
OWASP: A10:2021 - Server-Side Request Forgery (SSRF)
\end{lstlisting}
\end{small}

\textbf{Potongan Kode:}
\begin{small}
\begin{lstlisting}[style=promptstyle]
578  }
579  }
580  $cache_path .= md5($uri);
581
582  if ( ! $fp = @fopen($cache_path, 'w+b'))
583  {
584 >> log_message('error', 'Unable to write cache file: '.$cache_path);
585     return;
586  }
587
588  }
\end{lstlisting}
\end{small}

\textbf{Kontrol ISO/IEC 27001:2022:} A.8.28, A.8.29

\bigskip
\textbf{Catatan Tambahan:}

\textit{Penjelasan Kerentanan:}
The vulnerable code is located in the \texttt{\_write\_cache} method
within \texttt{Output.php}. The vulnerability lies in the way the
cache path is constructed using user input from the URI and query
string. The \texttt{\$cache\_path} is built by appending an MD5 hash
of \texttt{\$uri}, where \texttt{\$uri} may contain unsanitized
user-controlled data from \texttt{\$\_SERVER['QUERY\_STRING']}.
This allows an attacker to influence the file path used in
\texttt{fopen()}, potentially leading to unauthorized access to
server-side resources.

\textbf{Saran Remediasi:} \textit{Manual review required.}

\textbf{Penilaian Paket B:}

\begin{center}
\renewcommand{\arraystretch}{1.6}
\begin{tabular}{|l|p{7.5cm}|c|}
\hline
\textbf{Dimensi} & \textbf{Deskripsi} & \textbf{Skor (1--7)} \\
\hline
Akurasi & Seberapa akurat identifikasi kerentanan? & \\
\hline
Kejelasan & Seberapa jelas penjelasan yang diberikan? & \\
\hline
\textit{Actionability} & Seberapa \textit{actionable} rekomendasi perbaikan? & \\
\hline
Ketepatan ISO & Seberapa tepat pemetaan kontrol ISO 27001? & \\
\hline
\multicolumn{2}{|l|}{Komentar} & \\
\hline
\end{tabular}
\end{center}

% ============================================================
\section*{D.4 Catatan: Temuan FT-03 hingga FT-11}
% ============================================================

Temuan FT-03 hingga FT-11 seluruhnya merupakan kerentanan SSRF
(\textit{Server-Side Request Forgery}) yang bersumber dari pola
kode yang sama di \texttt{Output.php}, yaitu penggunaan
\texttt{\$uri} atau \texttt{\$filepath} yang berasal dari
data pengguna sebagai argumen operasi \textit{file system}.
Struktur instrumen untuk setiap temuan identik dengan FT-02,
dengan perbedaan hanya pada nomor baris (\textit{line number})
dan potongan kode yang ditampilkan. Penugasan kondisi B0/T1 ke
Paket A atau Paket B pada FT-03 hingga FT-11 diacak secara
independen menggunakan \texttt{random.seed(42)}.

Tabel berikut merangkum nomor baris dan penugasan paket
untuk setiap temuan FT-03 hingga FT-11:

\begin{center}
\renewcommand{\arraystretch}{1.3}
\begin{tabular}{|c|c|c|c|}
\hline
\textbf{Temuan} & \textbf{Baris} & \textbf{Paket A} & \textbf{Paket B} \\
\hline
FT-03 & 632 & B0 & T1 \\
\hline
FT-04 & 637 & T1 & B0 \\
\hline
FT-05 & 677 & T1 & B0 \\
\hline
FT-06 & 677 & T1 & B0 \\
\hline
FT-07 & 684 & T1 & B0 \\
\hline
FT-08 & 684 & B0 & T1 \\
\hline
FT-09 & 698 & B0 & T1 \\
\hline
FT-10 & 704 & B0 & T1 \\
\hline
FT-11 & 766 & T1 & B0 \\
\hline
\end{tabular}
\end{center}

\noindent\textit{Keterangan: B0 = keluaran \textit{pipeline} tanpa
\textit{AI Engine} (hanya data Semgrep mentah); T1 = keluaran
\textit{pipeline} dengan \textit{AI Engine} aktif (Qwen2.5-Coder 7b),
disertai penjelasan kerentanan dan saran remediasi.}
